%%%%%%%%%%  scan idx 659  |  folio 646  %%%%%%%%%%
\SGAPrintPage{646}%
\[
  \chi(F\otimes L)=\chi(F)
\]
\item
\emph{\(L\) is numerically equivalent to \(0\).}
\end{enumerate}
\emph{If \(X/k\) is projective and equipped with an ample invertible module \(H={O}_X(1)\), the preceding conditions are also equivalent to the following}:
\begin{enumerate}[label=(\roman*),start=4]
\item \emph{For every integer \(m\) (positive or negative), \(L^{\otimes m}(1)\) is ample.}
\item \emph{(If \(X\) is irreducible). For every pair of integers \(m,n\), we have}
\[
  \chi(L^{\otimes m}(n))=\chi({O}_X(n))
\]
\item \emph{(If \(X\) is irreducible of dimension \(r\ge2\)). We have}
\[
  \langle c_1(L)\cdot c_1(H)^{r-1}\rangle=0,\qquad
  \langle c_1(L)^2\cdot c_1(H)^{r-2}\rangle=0.
\]
\end{enumerate}

Let us prove that (i) implies (ii). Applying the d\'evissage lemma (EGA III 3.1) and 4.5 (i), we reduce to the case where \(X\) is integral and \(F={O}_X\). We must then prove that the polynomial \(\chi(L^{\otimes m})\) is constant. For this, we may replace \(L\) by a multiple, and thus assume that \(L\) is algebraically equivalent to \(0\). In this case, if \(L\) is deformed into \({O}\) by the algebraic family defined by an invertible sheaf \(M\) on some \(X\times_kT\) where \(T\) is connected, of finite type, the polynomial \(\chi(M_t^{\otimes m})\) does not depend on \(t\in T\) (EGA III 7.9.5); hence the implication.

Let \(Y\) be a closed integral curve in \(X\). Taking \(F={O}_Y\) in (ii), we find \(\chi({O}_Y)=\chi(L_Y)\). Thus the formula (``Riemann's theorem'') \(\chi(L_Y)=\langle c_1(L)\cdot Y\rangle+\chi({O}_Y)\) gives \(\langle c_1(L)\cdot Y\rangle=0\); hence (ii) implies (iii).

Let \(f:X'\to X\) be a Chow presentation: \(X'\) is projective and \(f\) is surjective. If we grant that (iii) implies (i) for \(X'\), then
%%%%%%%%%%  scan idx 660  |  folio 647  %%%%%%%%%%
\SGAPrintPage{647}%
by 4.5 the implication for \(X\) follows. Henceforth, then, suppose that \(X\) is projective, equipped with an ample sheaf \(H\).

Let \(X_i\) be the irreducible components of \(X\), equipped with the induced reduced structure. Applying 4.5 and EGA III 2.6.2 to the finite surjective morphism \(\coprod X_i\to X\) reduces the proof that (iii) implies (i) and (iv) to the case where \(X\) is integral.

If \(X\) is irreducible, it is clear that (ii) implies (v), that (iii) implies (vi), and, taking 2.1 into account, that (v) implies (vi), and that (vi) for \(X\) is equivalent to (vi) for \(X_{\red}\) (X 4.4.1).

It therefore remains to show that (vi) implies (i) and (iv), and that (iv) implies (i), assuming \(X\) integral. But then (vi) (resp. (iv)) implies that the family \(\mathcal L\) of classes of the \(L^{\otimes m}\) (\(m\in\mathbb N\)) satisfies 3.13 (iii) (resp. (iv)), and is therefore bounded. There thus exists an invertible sheaf \(L\) on some \(X_T\), with \(T\) of finite type over \(k\), which bounds \(\mathcal L\). Let \(p,q\) be two distinct integers such that \(L^{\otimes p}\) and \(L^{\otimes q}\) correspond to the same connected component of \(T\); then \(L^{\otimes(p-q)}\) is therefore algebraically equivalent to \(0\); hence (vi) (resp. (iv)) implies (i).

Finally, we likewise have that (vi) implies the existence of an integer \(N\) such that, for every integer \(m\), \(L^{\otimes m}\otimes H^{\otimes N}\) is ample. Then, for every \(m\), \((L^{\otimes m}\otimes H)^{\otimes N}\) is ample; hence (vi) implies (iv), which completes the proof.

\sgaStableTarget{sga6.e13.thm.4-7}\textbf{Theorem 4.7.} \emph{Let \(S\) be a noetherian scheme, and \(X\) a proper \(S\)-scheme. Then}:
\begin{enumerate}[label=(\roman*)]
\item \emph{The inclusion morphism \(\underline{\Pic}^\tau_{X/S}\to\underline{\Pic}_{X/S}\) is representable by an open immersion.}
\item
\emph{If \(X\) is projective and integral over \(S\), then \(\underline{\Pic}^\tau_{X/S}\to\underline{\Pic}_{X/S}\)
%%%%%%%%%%  scan idx 661  |  folio 648  %%%%%%%%%%
\SGAPrintPage{648}%
is also representable by a closed immersion.}
\item \emph{\(\underline{\Pic}^\tau_{X/S}\) is of finite type over \(S\), that is, the family of classes of invertible sheaves \(\tau\)-equivalent to \(0\) on the fibers of \(X/S\) is bounded.}
\end{enumerate}

Indeed, in (i) (resp. (ii)), we must treat the case of a morphism \(T\to\underline{\Pic}_{X/S}\), where \(T\) is an arbitrary \(S\)-scheme, and show that the inverse image of \(\underline{\Pic}^\tau_{X/S}\) is representable by an open (resp. and closed) subscheme of \(T\). It suffices to take \(T=\Spec(B)\) over an affine open \(\Spec(A)\) of \(S\). Then \(T\) is the limit of the filtering inverse system of affine \(S\)-schemes \(T_\alpha=\Spec(B_\alpha)\), where \(B_\alpha\) runs through the finite-type \(A\)-subalgebras of \(B\). By 3.1, the morphism in question comes from a morphism \(T_\alpha\to\underline{\Pic}_{X/S}\). We may therefore suppose that \(T\) is of finite type over \(S\).

The morphism \(T\to\underline{\Pic}_{X/S}\) corresponds to an f.p.p.f. morphism \(R\to T\) and an invertible sheaf \(L\) on \(X_R\). By f.p.p.f. descent, we may replace \(T\) by \(R\). After the base change \(R\to S\), taking into account the compatibility of the formation of \(\underline{\Pic}^\tau_{X/S}\) with every base change (4.1), we are thus reduced in (i) and (ii) to studying the case of a section \(S\to\underline{\Pic}_{X/S}\), defined by an invertible sheaf \(L\) on \(X\).

We must prove that the set \(S_0\) of \(s\in S\) such that \(L_s\) is \(\tau\)-equivalent to \(0\) is open, and also closed under the conditions of (ii); since \(S\) is noetherian, this means three things: for an \(s\in S_0\),
\begin{enumerate}[label=(\alph*)]
\item every generalization of \(s\) remains in \(S_0\).
\item every specialization of \(s\), near \(s\), remains in \(S_0\).
\item under the conditions of (ii), every specialization of \(s\) remains in \(S_0\).
\end{enumerate}

%%%%%%%%%%  scan idx 662  |  folio 649  %%%%%%%%%%
\SGAPrintPage{649}%
Let \(f:X'\to X\) be a Chow presentation: \(X'\) is projective over \(S\); and \(f\) is surjective. If we grant (a), (b), (c) and (iii) for \(X'\), then they follow for \(X\) by 4.3 and 3.5. We may therefore suppose that \(X\) is projective over \(S\).

Let \(H\) be ample for \(X/S\). If, for an \(s\in S\), \((L^{\otimes m}\otimes H)_s\) is ample, then \((L^{\otimes m}\otimes H)_t\) is ample for every generalization \(t\) of \(s\) (EGA III 4.7.1). By the equivalence of 4.6 (i) and (iv), we have therefore established (a).

For (b), (c) and 4.7 (iii), we may suppose that \(S\) is integral, with \(s\) the generic point in (b) and (c) (For (b) and (c), this is clear, and for (iii), it follows from 3.4 (i) applied to the morphism \(\coprod S_i\to S\), where the \(S_i\) are the irreducible components of \(S\), equipped with the induced reduced structure).

There exist an integral scheme \(S'\) and an f.p.p.f. morphism \(S'\to S\) such that every irreducible component \(X'_i\) of \(X_{S'}\), equipped with the induced reduced structure, is integral over \(S'\) (XII 1.3). To prove (b) and (iii), it suffices to do so for \(X_{S'}/S'\), (in (b), by descent; in (iii), by noetherian induction and 3.4 (i)).

Finally, consider the canonical surjective morphism \(\coprod X'_i\to X_{S'}\). Then (b) and (iii) for \(X_{S'}/S'\) follow from (c) and (iii) for the \(X'_i/S'\), by 4.3, 3.5 and 3.7.

We have therefore reduced to proving (c) and (iii), assuming \(S\) integral with generic point \(s\), and \(X\) projective and integral over \(S\). Let \({O}_X(1)\) be a very ample sheaf for \(X/S\). Then, for every \(m,n\), \(\chi(L_t^{\otimes m}(n))\) and \(\chi({O}_{X_t}(n))\) do not depend on \(t\in S\) (EGA III 7.9.4). Thus (c) follows from 4.6 (i) \(\Longleftrightarrow\) (v), which also implies that the family of classes of invertible sheaves \(\tau\)-equivalent to \(0\) on the fibers of \(X/S\) has just
%%%%%%%%%%  scan idx 663  |  folio 650  %%%%%%%%%%
\SGAPrintPage{650}%
one Hilbert polynomial; thus (iii) follows from 2.11.

\sgaStableTarget{sga6.e13.rem.4-8}\textbf{Remark 4.8.} Let \(X\) be proper and integral over noetherian \(S\). Without a projectivity hypothesis on \(X/S\), it is unknown whether it remains true that \(\underline{\Pic}^\tau_{X/S}\to\underline{\Pic}_{X/S}\) is a \uline{closed} immersion, except when \(X\) is normal over \(S\) and \(\underline{\Pic}_{X/S}\) is \uline{representable} (FGA 236, 2.3).

% Added in pages 654--669
\providecommand{\NS}{\operatorname{NS}}
\providecommand{\Picard}{\operatorname{Pic}}
\providecommand{\Et}{\mathrm{et}}
\providecommand{\rk}{\operatorname{rk}}
\providecommand{\im}{\operatorname{Im}}
\providecommand{\HH}{\mathrm H}


\sgaStableTarget{sga6.e13.sec.5}\subsection*{5. Finiteness theorems of ``N\'eron-S\'everi'' type.}

\sgaStableTarget{sga6.e13.thm.5-1}\textbf{Theorem 5.1.} \emph{Let \(X\) be proper over noetherian \(S\); then the N\'eron-S\'everi groups \(\underline{\Pic}_{X_K/K}(K)/\underline{\Pic}^0_{X_K/K}(K)\) of the geometric fibers \(X_K/K\) of \(X/S\) are finitely generated, and their ranks and the orders of their torsion subgroups \(\underline{\Pic}^\tau_{X_K/K}(K)/\underline{\Pic}^0_{X_K/K}(K)\) remain bounded.}

Indeed, by noetherian induction, we may restrict ourselves to proving 5.1 with \(S\) replaced by an arbitrary nonempty open of \(S_{\red}\); thus we may suppose that \(\underline{\Pic}^\tau_{X/S}\) is representable by an \(S\)-scheme of finite type (XII 1.2 and 4.7 (iii)). Then the order of \(\underline{\Pic}^\tau_{X_K/K}(K)/\underline{\Pic}^0_{X_K/K}(K)\) is equal to the number of connected components of \((\Pic^\tau_{X/S})_K\), a number that remains the same over a nonempty open of \(S\) (EGA IV 9.7.9).

It remains to bound the rank of \(\underline{\Pic}_{X_K/K}(K)/\underline{\Pic}^\tau_{X_K/K}(K)\). As in the proof of 4.7 (iii), we readily reduce to the case where \(S\) is integral, and \(X\) is projective and integral over \(S\). Let \(r\) be the common dimension of the fibers of \(X/S\). If \(r\) is \(0\) (resp. \(1\)), then this rank is clearly \(0\) (resp. \(1\)).

If \(r\) is \(\ge2\), we shall reduce to the case where \(r=2\) and \(X\) is smooth over \(S\). After restricting \(S\), we can find a linear section \(Y\) of \(X\) which is
